\section{Predicted bands and the refined tables}\label{sec:bands}

Starting from the fixed seed and the table $d_0=w_T$, we now define
$d_1,\ldots,d_S$. The construction proceeds by node height: when defining
the value at a parent and index $t$, all child tables through index $t$
are already defined. This is a mathematical ordering, not a requirement
to evaluate those tables in full. The algorithm will request only the
entries needed by the choices below.

There are two ways to choose a current-index action. A shared sample
compares rounded actions at nearby representatives; extending the chosen
cost by a cone gives a spatially regular estimate. An online rule uses
earlier child values to choose an action with small expected accumulated cost
relative to every fixed action in a predicted band. The parent update
combines these two estimates. We first construct the band and its child
cost bounds, then specify the online rule and the shared samples.

\subsection{A band predicted from prefix costs}

At a representative $(b,r)$ of a nonleaf, define
\[
 Z=\min_{\alpha\in\mathcal A_{b,r}}V_r(\alpha;m).
\]
Skip the band computations at this representative unless $Z$ is finite
and strictly positive.  Otherwise set $\widehat v=z(I,r)/Z$, the net
shift per unit of predicted cost. If shift accumulated proportionally to
this cost along an alignment, a prefix of cost $C$ would end at shift
$p_0+\widehat v C$. We use this relation to predict the shift at each
child cut. We call $|\widehat v|\le1-\tau$ the
\emph{velocity test}, and skip unless it holds.

For an internal cut $1\le j<M$ and a real shift $p$, let
$\Phi_j(p)$ be the minimum prefix cost using the first $j$ child
restrictions of the wide action set.  More explicitly, minimize
\[
 |s_1^--r^-|+\sum_{i=1}^j m(I_i,s_i)
 +\sum_{i<j}|s_{i+1}^--s_i^+|+|s_j^+-(l_j+p)|
\]
over such prefix states.  The wide action family is nonempty whenever
the band reaches this step, so these prefix minima are finite.  Each
function in the minimum is 1-Lipschitz in $p$, hence so is $\Phi_j$.

The prefix cost depends on its proposed ending shift. The predicted
shift therefore solves the approximate fixed-point relation
$p=p_0+\widehat v\Phi_j(p)$. The velocity test keeps the Lipschitz
constant of $\widehat v\Phi_j$ below one, so all sufficiently accurate
solutions lie in a short interval. The band below imposes this relation
at every cut.

Let $\mathcal C_{b,r}$ be the following set of shared-endpoint actions
on the scale-$b$ grid.  The outside endpoints are exactly $r$,
consecutive child endpoints agree, all endpoint shifts are within
$6b$ of $p_0=r^--l$, and the shift $p$ at every internal cut satisfies
\begin{equation}\label{eq:band}
 |p-p_0-\widehat v\Phi_j(p)|\le\tau\theta b/M.
\end{equation}
All child target states must be valid, so these shared endpoints are in
nondecreasing order.  Skip if $\mathcal C_{b,r}$ is empty.

The possible shifts at any fixed cut have diameter at most
$2\theta b/M$.  In fact two allowed shifts obey
\[
 |p-p'|\le2\tau\theta b/M+
 |\widehat v|\,|\Phi_j(p)-\Phi_j(p')|
 \le2\tau\theta b/M+(1-\tau)|p-p'|.
\]
No convexity of the allowed shift set is needed.

We need a uniform cost bound at each child for both action selection and
rounding. For a nonempty band, let $h_i$ be the smallest power of two at least
\[
 \max\left\{b/M,\max_{\alpha\in\mathcal C_{b,r}}
                  w_0(I_i,s_i(\alpha))\right\},
 \qquad h_\Sigma=\sum_{i=1}^M h_i.
\]
Powers of two below one are allowed here. Every band state at child $i$
has $w_0$ value at most $h_i$. These bounds will normalize the online
costs and determine a common rounding spacing for each child scale.
Skip if $h_\Sigma>64Fb$.
Choose a canonical band action in a fixed ordering and round its two
coordinates separately at child $i$ down to multiples of
$\max\{1,\theta h_i\}$.  Denote the resulting action by
$\beta(b,r)$. Different children may round a shared endpoint to different
positions, so $\beta$ can have gaps or overlaps. It is used as a wide
action with its connection charge, rather than as a band action.
Finally, skip unless it belongs to $\mathcal A_{b,r}$
and
\begin{equation}\label{eq:rounded-guard}
 |z(I_i,s_i(\beta))|\le2h_i\qquad(1\le i\le M).
\end{equation}
Representatives passing all these tests will be called unskipped.
These tests affect only the band and rounded-action computations;
they never remove an action or representative from
$\mathcal B_I$.

Under the fixed good seed, the last two guards on $\beta$ pass
whenever the $h_\Sigma$ test passes.  Indeed,
$\theta h_i\ge\theta b/M\ge b/P$.  All relevant quantities are
powers of two, so the rounded spacing is a multiple of $g_b$.
Rounding preserves the order of each child's two endpoints and the range
$[0,|y|]$. If its
spacing is one, its error is zero; otherwise each coordinate error is
less than $\theta h_i$.  Since $|z|\le w_0\le h_i$ before rounding,
\[
 |z(I_i,s_i(\beta))|\le h_i+2\theta h_i\le2h_i.
\]
The endpoint shift radius after rounding is at most
$6b+64F\theta b<10b$.  Section~\ref{sec:comparison} will prove the
earlier $h_\Sigma$ test on the needed comparison nodes.

\subsection{The finite-mixture interface for online choices}

The online rule uses child values at indices less than $t$ to select
one band action at time $t$. Its accumulated current-index cost will be
compared with that of each fixed band action in Section~\ref{sec:regret}.
We encode a mixture of actions by its child-state marginals, with weights
chosen to recover the unnormalized sum of child costs.

Fix one unskipped band $\mathcal C=\mathcal C_{b,r}$. Its coordinates
are pairs $(i,s)$, one for each child and each state appearing at that
child in a band action.  Let $d_{\mathcal C}$ be their number.  Map
each action $\alpha$ to the probability vector $\rho(\alpha)$ that
places mass $h_i/h_\Sigma$ at $(i,s_i(\alpha))$ for every child,
and zero at its other coordinates.  Let
$\mathcal Q=\operatorname{conv}\{\rho(\alpha):\alpha\in\mathcal C\}$.
Every vector in this convex hull has the same total mass $h_i/h_\Sigma$
on the coordinates of child $i$.

Define coordinate vectors
\begin{align*}
 \widehat g_j(i,s)&=-d_{j-1}(I_i,s)/h_i,\\
 \Delta_j(i,s)&=(d_{j-1}(I_i,s)-d_j(I_i,s))/h_i,\\
 u_j&=\widehat g_j+(1-2\kappa)\Delta_j.
\end{align*}
Here $\widehat g_j$ predicts the negative normalized current cost from
the preceding index, and $\Delta_j$ is the observed correction:
\[
 \widehat g_j(i,s)+\Delta_j(i,s)=-d_j(I_i,s)/h_i,
 \qquad
 -h_\Sigma\langle\widehat g_j,\rho(\alpha)\rangle
 =\sum_i d_{j-1}(I_i,s_i(\alpha)).
\]
The fixed mass $h_i/h_\Sigma$ in each child block cancels the
normalization by $h_i$. The objective below uses $\Delta_j$ only for
$j<t$ and uses $\widehat g_t$, so its required child indices are all
strictly earlier than $t$.

With entropy $\mathcal H(\rho)=-\sum_k\rho_k\ln\rho_k$, using
$0\ln0=0$, let $\rho_t$ maximize
\begin{equation}\label{eq:online-objective}
 \mathcal H(\rho)+
 \eta\left\langle\sum_{j<t}u_j+\widehat g_t,\rho\right\rangle
 \quad\hbox{over }\rho\in\mathcal Q.
\end{equation}
The algorithm computes a finite mixture of band actions whose vector
of expected weights is within $H^{-3}$ in $\ell_1$ distance of
$\rho_t$, then samples that mixture exactly.  Section~\ref{sec:numerics}
constructs this mixture using rational arithmetic and proves its stated
accuracy and sampling cost.  Thus~\eqref{eq:online-objective} specifies
an input-output contract for a finite algorithm, not an optimization
oracle.

Only the child histories are needed for this rule, not a transcript of
earlier queries at the parent.  Parent draws may be keyed by
$(I,q,b,t)$ in addition to pass and copy identifiers.  Group choices
and table values are cached on their complete keys, and every first
parent draw uses a fresh source independent of the complete child
tables.  Section~\ref{sec:work} makes the resulting deferred evaluation
and its costs explicit.

\subsection{Shared group samples and the update}

For each node $I$, scale $b$, cell from
Lemma~\ref{lem:cone-locality}, and time $1\le t\le S$, group all
unskipped representatives in that whole cell.  The following procedure
chooses one member of each nonempty group.  It is defined from complete
child tables; later we evaluate its required entries on demand.

Sampling and rounding work together here. Sharing the child draws
across a cell lets all its representatives reuse the same queried child
labels. Rounding then limits how many different states at one such child
can occur across the group. This combination, rather than band width
alone, bounds the number of current-index recursive calls.

For each power of two $h$ in $[b/M,64Fb]$, set
\[
 k_h=\min\{M,\ceil{Q_1Mh/b}\}.
\]
For each member $(b,r)$, we estimate the child sum in
$V_r(\beta(b,r);d_t)$ by separating the children according to their
$h_i$.  If $k_h=M$, include the contribution with $h_i=h$ exactly.
Otherwise draw $k_h$ independent uniform child indices with replacement,
using the same draws for every member of this group at this scale.
For each sampled child include its $d_t$ value if its $h_i=h$, and zero
otherwise; multiply the sampled sum by $M/k_h$.  Add the connection
charge exactly.  Select a member minimizing this estimate, breaking
ties by a fixed ordering.  Then evaluate its actual action cost using
all children.  Samples used for different parent groups and times are
independent and are independent of descendant randomness.

For example, when $k_h<M$, there are $O(Q_1Mh/b)$ sampled child
labels, and at most $O(\theta^{-2}(1+b/h))$ rounded states at each
label across the cell. Their product is $M\ell^{O(1)}$.
Lemma~\ref{lem:current-index-budget} proves this count and the analogous
bound for exact scales, including executions with inaccurate seeds.

Let $\Pi_t(I,q)$ be the minimum of
$\mathcal B_I(d_{t-1})(q)$ and the actual selected group costs, each
extended from its selected representative by $L|q-r|_1$.  Compute this
minimum by~\eqref{eq:local-scales} and whole-cell enumeration.  A
group's choice is cached for the whole cell, independent of which
query first requests it.  Lemma~\ref{lem:cone-locality} shows that this
local computation equals the global cone minimum for a good seed.

There is also an online action query at each scale passing
\eqref{eq:local-scales}.  Put $r=r_b(q)$ and proceed only if this is
an unskipped representative.  Draw one action
$\alpha_t\in\mathcal C_{b,r}$ by the rule just defined,
which uses child histories with index less than $t$.  Evaluate
\[
 L|q-r|_1+\sum_i d_t(I_i,s_i(\alpha_t)).
\]
Let $Y_t(I,q)$ be the minimum of these online query values, or
$+\infty$ if there are none.  At a nonleaf define
\begin{equation}\label{eq:refined-update}
 d_t(I,q)=
 \begin{cases}
 0,&w_0(I,q)=0,\\[2pt]
 \displaystyle\min\left\{\Pi_t(I,q),
       \max\left\{\frac{\Pi_t(I,q)}{1+\delta},Y_t(I,q)\right\}\right\},
       &w_0(I,q)>0.
 \end{cases}
\end{equation}
Leaves continue to use their exact values. At a positive $w_0$ state,
this update truncates $Y_t$ to the interval
$[\Pi_t/(1+\delta),\Pi_t]$: an inexpensive online action can improve
the value only within this interval. The upper endpoint retains the
Bellman comparison, while the lower endpoint transfers the spatial
regularity of $\Pi_t$ to $d_t$ up to a factor $1+\delta$.
For a fixed pass, these definitions are well founded by induction on
node height: current-index values are requested only at children, and
at a fixed node the remaining refinement values have earlier indices.
Section~\ref{sec:work} gives the complete lazy evaluation order across passes.

\subsection{Deterministic invariants}

\begin{lemma}\label{lem:refinement-invariants}
For a seed satisfying Theorem~\ref{thm:refinement}, every outcome of
the refinement obeys
\begin{align}
 \max\{c/a,|z|\}&\le d_t\le d_{t-1}\le w_0,
 \label{eq:refinement-invariants}\\
 d_t(I,q)&\le(1+\delta)d_t(I,q')+L|q-q'|_1.
 \label{eq:refinement-lipschitz}
\end{align}
At nonleaves one also has
\begin{equation}\label{eq:bellman-sandwich}
 \mathcal B_I(d_t)(q)\le d_t(I,q)
           \le\mathcal B_I(d_{t-1})(q).
\end{equation}
For $d_0$, the spatial bound holds with Lipschitz constant $L$ and
without the factor $1+\delta$.
\end{lemma}

\begin{proof}
Induct on node height, and at a fixed node on $t$.  Every extra
candidate in $\Pi_t$ is a current-value wide-action cone.  The Bellman
candidate uses earlier child values, which are at least the current
ones by induction.  Hence $\Pi_t\ge\mathcal B_I(d_t)$, whereas
$\Pi_t\le\mathcal B_I(d_{t-1})$.  Each online candidate is a
current-value wide-action cone too, so
$Y_t\ge\mathcal B_I(d_t)$, including the empty case.  Taking the
specified minimum and maximum in~\eqref{eq:refined-update} proves
\eqref{eq:bellman-sandwich}.  In particular, no lower bound is inferred
from $\Pi_t/(1+\delta)$ alone.

At $t=1$, \eqref{eq:warmup-invariants} now gives $d_1\le d_0$.
At later times the left inequality in~\eqref{eq:bellman-sandwich}
at time $t-1$ gives $d_t\le d_{t-1}$.  The lower bound follows from
Bellman preservation and the child induction.

The global cone interpretation makes $\Pi_t$ $L$-Lipschitz, whether
or not any group estimate is accurate.  Moreover
$\Pi_t/(1+\delta)\le d_t\le\Pi_t$.  Thus
\[
 d_t(I,q)\le\Pi_t(I,q)
 \le\Pi_t(I,q')+L|q-q'|_1
 \le(1+\delta)d_t(I,q')+L|q-q'|_1.
\]
If $w_0(I,q)=0$, the Bellman and $\Pi_t$ values are zero, and the
explicit zero branch agrees with these inequalities.  The assertions
for $d_0$ follow from the warmup.
\end{proof}

For execution on an arbitrary seed, the local construction still gives
nonnegative values at most its local $w_0$.  All representative lists
and skipped-band tests remain finite.  We use these unconditional
facts for work bounds, and do not extend the correctness invariants of
Lemma~\ref{lem:refinement-invariants} to inaccurate seeds.

\subsection{Accuracy of one group selection}

For a band action define
$x_t(\alpha)=\sum_i d_t(I_i,s_i(\alpha))$.
Two band actions differ by at most $4\theta b/M$ in each child's
endpoint-pair metric.  Rounding the canonical action adds at most
$2\theta h_i$, so
$|s_i(\alpha)-s_i(\beta)|_1\le6\theta h_i$.
Also $\conn_r(\beta)\le2\theta h_\Sigma$, because the unrounded
action has shared endpoints.  It follows that
\begin{equation}\label{eq:rounded-action}
 \begin{split}
 V_r(\beta;d_t)&\le(1+\delta)x_t(\alpha)
                    +(6L+2)\theta h_\Sigma,\\
 d_t(I_i,s_i(\beta))&\le2h_i.
 \end{split}
\end{equation}
The first inequality uses~\eqref{eq:refinement-lipschitz}.  For the
second, use $d_t\le w_0$ and the $L_0$-Lipschitz bound to obtain
$w_0(I_i,s_i(\beta))\le h_i+2L_0\theta h_i\le2h_i$.

\begin{lemma}[One-group accuracy]\label{lem:group-accuracy}
Under the seed hypothesis of Theorem~\ref{thm:refinement}, fix complete
child tables and one group at time $t$.  With conditional
probability at least $1-H^{-4}$, all action-cost estimates used to choose
its representative are within $\gamma b$ of their actual costs.
Consequently, for any specified unskipped member $(b,r)$ of this
group, on the same event, simultaneously for every valid parent state
$q$ and every $\alpha\in\mathcal C_{b,r}$,
\begin{equation}\label{eq:group-accuracy}
 \Pi_t(I,q)\le L|q-r|_1+(1+\delta)x_t(\alpha)+e_*b,
 \qquad e_*=2\gamma+512LF\theta\le\kappa\delta.
\end{equation}
Here the inequality uses the global cone minimum, equivalently its
local evaluation from Lemma~\ref{lem:cone-locality}.
\end{lemma}

\begin{proof}
Conditioning fixes the bands, the $h_i$, the rounded actions, and all
their child costs.  The parent's group samples remain independent of
these data.  For a fixed member and sampled scale $h$, normalize one
sample to
\[
 X=\frac Mb\ind_{h_i=h}\,d_t(I_i,s_i(\beta)),
\]
where $i$ is uniform in $\{1,\ldots,M\}$.  By
\eqref{eq:rounded-action},
\[
 0\le X\le2Mh/b,
 \qquad \E X\le2h_\Sigma/b\le128F,
 \qquad \operatorname{Var}(X)\le256FMh/b.
\]
When $k_h<M$, the number of independent samples is at least
$Q_1Mh/b$.  Bernstein's inequality
\cite[Theorem~4.16]{Tropp2021Probability}, with normalized error
$\alpha_0=\gamma/(2\ell)$, gives a failure probability at most
\[
 2\exp\left(-\frac{Q_1\alpha_0^2}{512F+4\alpha_0/3}\right)
 \le2\exp(-\ell^{37}/3000).
\]
When $k_h=M$, the contribution is exact and has no failure event.
There are at most $2\ell$ scales and at most $(P+1)^2$ group members.
A union bound is at most $H^{-4}$ at the absolute cutoff.  Summing
the errors over scales gives error at most $\gamma b$ for every
member.  Shared samples between members cause no difficulty for this
union bound.

On this event, estimated-cost minimization followed by actual
evaluation selects a cost at most that of the specified member plus
$2\gamma b$.  The selected representative is in the same cell, so
its endpoint-pair distance from $r$ is at most $2\theta b$.
Using~\eqref{eq:rounded-action} and $h_\Sigma\le64Fb$ therefore
gives~\eqref{eq:group-accuracy}, since its additive coefficient is
at most
\[
 2\gamma+(384LF+128F+2L)\theta
 \le2\gamma+512LF\theta\le\kappa\delta.
\]
The comparison with every band action is simultaneous because
\eqref{eq:rounded-action} is deterministic and uniform in that
action.
\end{proof}

This is a conditional statement for one whole cell and time.  The
accuracy proof will not require all group samples throughout the
recursion to succeed together.  Instead it will charge each group's
failure probability inside an expected regret bound at its comparison
node.
